\section{Explicit Feynman coefficients for 6D hCS observables (Wave-13 H3/20)}
\label{wn:sec:wave13-h3-feynman-6d-hCS}

\providecommand{\ClaimStatusTheorem}{\textsc{[theorem]}}
\providecommand{\ClaimStatusConjectured}{\textsc{[conjectural]}}
\providecommand{\ClaimStatusHeuristic}{\textsc{[heuristic]}}
\providecommand{\ClaimStatusDefinition}{\textsc{[definition]}}
\providecommand{\kch}{\kappa_{\mathrm{ch}}}
\providecommand{\kcat}{\kappa_{\mathrm{cat}}}
\providecommand{\kBKM}{\kappa_{\mathrm{BKM}}}
\providecommand{\kfib}{\kappa_{\mathrm{fiber}}}
\providecommand{\cA}{\mathcal{A}}
\providecommand{\cO}{\mathcal{O}}
\providecommand{\Obs}{\mathrm{Obs}}
\providecommand{\hCS}{\mathrm{hCS}}
\providecommand{\CC}{\mathbb{C}}
\providecommand{\RR}{\mathbb{R}}
\providecommand{\gfrak}{\mathfrak{g}}
\providecommand{\uone}{\mathfrak{u}(1)}
\providecommand{\sun}{\mathfrak{su}(N)}
\providecommand{\Cas}{C_2(\gfrak)}
\providecommand{\Tr}{\mathrm{tr}}

\subsection{Target and scope}

Building on E5, which supplied Feynman rules and superficial power-counting, H3
traces every step of the lowest-order Feynman coefficients for 6D hCS on
$X = \CC^3$ with gauge algebra $\gfrak$ and action
\[
S_{\mathrm{cl}}[\cA] \;=\; \int_{\CC^3}\Omega \wedge \Tr\!\left(
\tfrac{1}{2}\,\cA\,\bar\partial\cA + \tfrac{1}{3}\,\cA^3\right),
\qquad \Omega = dz_1 dz_2 dz_3,
\]
writing each observable coefficient symbolically with the Bochner--Martinelli
propagator explicit.

\subsection{Cycle 1 --- Two-point function $\langle A_{0,1}(z_1) A_{0,1}(z_2)\rangle$}

\ClaimStatusDefinition\ In free ($\gfrak=\uone$, or order-$\hbar^0$ non-abelian)
theory the two-point function of $\cA \in \Omega^{0,1}(\CC^3)\otimes\gfrak$ is
the Bochner--Martinelli propagator. Let $\phi_1,\phi_2 \in
\Omega^{0,1}_c(\CC^3)\otimes\gfrak$ be compactly supported test 1-forms. Then
\[
\bigl\langle \cA^a(\phi_1)\,\cA^b(\phi_2) \bigr\rangle_0
\;=\; \hbar\,\delta^{ab}\int_{\CC^3\times\CC^3}\Omega_z\Omega_w\,
\phi_1(z)\wedge G(z,w)\wedge \phi_2(w),
\]
with
\[
G(z,w) \;=\; \frac{2!}{(2\pi i)^3}\,
\sum_{k=1}^{3}(-1)^{k-1}\,\frac{\overline{(z_k{-}w_k)}}{\|z-w\|^{6}}\,
\widehat{d\bar z_k}\wedge dw_1 dw_2 dw_3,
\]
the closed $(0,2)$-form in $z$, holomorphic top-form in $w$, away from the
diagonal $\Delta\subset\CC^3\times\CC^3$. Here $\widehat{d\bar z_k}$ denotes
$d\bar z_j\wedge d\bar z_l$ with $\{j,k,l\}=\{1,2,3\}$ in cyclic order.

\ClaimStatusTheorem\ \emph{Disjoint-support evaluation.} For
$\mathrm{supp}(\phi_1)\cap\mathrm{supp}(\phi_2)=\emptyset$, $G(z,w)$ is smooth on
the support of the product, the integral converges absolutely, and
$\bigl\langle \cA^a(\phi_1)\,\cA^b(\phi_2)\bigr\rangle_0$ equals the pairing
$\hbar\delta^{ab}\cdot(\phi_1,G\phi_2)_{L^2}$ with no regularisation needed.
The short-distance expansion as $z\to w$ is $G\sim C_3/r^5$ where
$r = \|z-w\|$ and $C_3 = 2!/(2\pi i)^3$, consistent with the $(2n{-}1)=5$
Bochner--Martinelli singularity in $n=3$.

\subsection{Cycle 2 --- Tree-level three-point function}

\ClaimStatusTheorem\ \emph{Abelian case.} The action
$S_{\mathrm{cl}}^{\uone}=\tfrac{1}{2}\int\Omega\wedge\cA\bar\partial\cA$ is
purely quadratic; Wick contractions generate only Gaussian moments, hence
\[
\bigl\langle \cA(\phi_1)\cA(\phi_2)\cA(\phi_3)\bigr\rangle_{\uone}^{\mathrm{tree}}
\;=\; 0.
\]
At all orders in $\hbar$, all odd correlators vanish and all even correlators
are products of two-point functions.

\ClaimStatusTheorem\ \emph{Non-abelian tree-level.} The cubic vertex
$\tfrac{1}{3}\Tr(\cA^3)$ supplies
\[
V_{abc}(w) \;=\; f_{abc}\,\delta^{(6)}(w\text{-insertion})\qquad
\text{(local trivalent vertex, symmetric in $abc$)}.
\]
The connected tree-level three-point correlator is
\[
\begin{aligned}
\bigl\langle \cA^{a_1}(\phi_1)\cA^{a_2}(\phi_2)\cA^{a_3}(\phi_3)
\bigr\rangle_0^{\mathrm{tree,conn}}
&= \hbar^2\,f^{a_1 a_2 a_3}
\int_{(\CC^3)^4}\!\!\!\!\Omega_{z_1}\Omega_{z_2}\Omega_{z_3}\Omega_w\\
&\quad\times\,\phi_1(z_1)\wedge G(z_1,w)\wedge
\phi_2(z_2)\wedge G(z_2,w)\wedge\\
&\quad\times\, \phi_3(z_3)\wedge G(z_3,w),
\end{aligned}
\]
i.e.\ a single bulk vertex at $w$ with three external legs convolved against the
propagators. The diagram is:
\[
\begin{array}{c}
\text{\texttt{%
      phi1 ---G---\Big[w\Big]---G--- phi2%
}}\\[-2pt]
\text{\texttt{%
                       |%
}}\\[-2pt]
\text{\texttt{%
                       G%
}}\\[-2pt]
\text{\texttt{%
                       |%
}}\\[-2pt]
\text{\texttt{%
                     phi3%
}}
\end{array}
\]

\subsection{Cycle 3 --- One-loop bubble diagram}

\ClaimStatusTheorem\ \emph{Abelian case.} No trivalent vertex exists; the
bubble vanishes: $I_2^{(1),\uone} = 0$. Equivalently, the $\uone$ theory is
exactly Gaussian and $Z_{\uone}[\phi] = \exp\bigl(\tfrac{1}{2\hbar}\int\Omega
\wedge\phi\bar\partial\phi\bigr)$ is $\hbar$-exact.

\ClaimStatusTheorem\ \emph{Non-abelian bubble.} Two trivalent vertices connected
by two internal propagators, two external legs $\phi_1,\phi_2$:
\[
\begin{array}{c}
\text{\texttt{%
      phi1 ---\Big[V_1\Big]======\Big[V_2\Big]--- phi2%
}}\\[-2pt]
\text{\small (two internal $G$'s between $V_1,V_2$; one external $G$ on each end)}
\end{array}
\]
Applying Feynman rules, with symmetry factor $1/2$ from the internal-propagator
exchange,
\[
\begin{aligned}
I_2^{(1)}[\phi_1^{a_1},\phi_2^{a_2}]
&= \frac{\hbar^2}{2}\,f^{a_1 b c}\,f^{a_2}{}_{bc}
\int_{(\CC^3)^4}\!\!\Omega_{z_1}\Omega_{z_2}\Omega_{w_1}\Omega_{w_2}\\
&\quad\times\phi_1(z_1)\wedge G(z_1,w_1)\wedge
G(w_1,w_2)\wedge G(w_2,w_1)\wedge G(w_2,z_2)\wedge\phi_2(z_2).
\end{aligned}
\]
The colour factor collapses via $\sum_{bc}f^{a_1 bc}f^{a_2}{}_{bc}
= \Cas\,\delta^{a_1 a_2}$ (definition of quadratic Casimir in the adjoint),
so
\[
I_2^{(1)} \;=\; \frac{\hbar^2\,\Cas}{2}\,\delta^{a_1 a_2}
\int_{\CC^3\times\CC^3}\!\!\Omega_{w_1}\Omega_{w_2}\,
[\phi_1\!\cdot\!G](w_1)\,\bigl[G(w_1,w_2)\wedge G(w_2,w_1)\bigr]\,
[\phi_2\!\cdot\!G](w_2),
\]
where $[\phi\cdot G](w)=\int\Omega_z\,\phi(z)\wedge G(z,w)$ is the propagator
convolution. For $\gfrak=\sun$ the standard normalisation gives $\Cas = 2N$.

\subsection{Cycle 4 --- Divergence structure of $\int G^2$}

\ClaimStatusTheorem\ \emph{Short-distance behaviour.} Setting $r = \|w_1-w_2\|$
and extracting the radial singularity, each $G$-factor contributes $r^{-5}$
(the $(2n{-}1)=5$ Bochner--Martinelli degree at $n=3$). The product
$G(w_1,w_2)\wedge G(w_2,w_1)$ behaves as $C_3^2\,r^{-10}$ (times angular
top-form). The integration measure along the diagonal in real coordinates
$w_2 - w_1 = r\cdot(\text{angular})$ over $\CC^3\cong\RR^6$ contributes
$r^5\,dr$. Hence
\[
\int_0^\Lambda^{-1}\!r^5\cdot r^{-10}\,dr
\;=\; \int_0^\Lambda^{-1}\!r^{-5}\,dr \;\sim\; \Lambda^{4}
\qquad (\text{UV divergent as }\Lambda\to\infty).
\]

\ClaimStatusHeuristic\ \emph{Reconciling with E5 power-counting.} The naive
BV superficial divergence $\omega = 6\cdot 1 - 4\cdot 2 = -2$ used
\emph{effective} propagator weight $-4$ (not $-5$), appropriate for the
$\bar\partial$-heat-kernel regularisation $G^\epsilon_L = \int_\epsilon^L
ds\,K_s$ which softens the short-distance behaviour. The raw
Bochner--Martinelli kernel has weight $-5$ in 6D; the heat-kernel cut-off
converts this to $\omega_{\mathrm{eff}}=-2$, hence logarithmic rather than
power divergence in the regulated theory. Both accountings agree: the
unregulated $r^{-10}$ distribution is not well-defined on $\Delta$, and the
minimal scheme produces a $\log\Lambda$ residue after BV renormalisation.

\subsection{Cycle 5 --- Wilsonian counterterm at one loop}

\ClaimStatusDefinition\ \emph{Regularised propagator.} Introduce the Costello
heat-kernel propagator
\[
G_{\epsilon}^{L}(z,w) \;=\; \int_\epsilon^L ds\,K_s^{\bar\partial}(z,w),\qquad
K_s^{\bar\partial}(z,w) \;=\; \frac{1}{(4\pi s)^3}\,e^{-\|z-w\|^2/(4s)}\,
(\bar\partial^*\text{-exact companion}),
\]
with $\epsilon\to 0$ UV regulator and $L\to\infty$ IR regulator.

\ClaimStatusTheorem\ \emph{One-loop counterterm.} After heat-kernel
regularisation, the bubble integrand $G_\epsilon^L\wedge G_\epsilon^L$ produces
a short-distance residue
\[
\int_0^L\!ds_1\int_0^L\!ds_2\,\frac{1}{(s_1+s_2)^3}
\;=\; \log(L/\epsilon) + \text{(finite)},
\]
(after reducing the 6D heat-kernel integral to the scale variable
$s=s_1+s_2$ via Fubini). The minimal-subtraction BV counterterm is therefore
\[
S_{\mathrm{c.t.}}^{(1)}[\cA] \;=\; -\,\hbar\,\frac{\Cas}{(4\pi)^3}\,\log(L/\epsilon)
\int_{\CC^3}\Omega\wedge \Tr(\cA\,\bar\partial\cA) + \cO(\hbar^2),
\]
a wave-function-renormalisation counterterm proportional to $\log\Lambda$
times the free kinetic operator. After adding $S_{\mathrm{c.t.}}^{(1)}$ the
renormalised effective action is finite as $\epsilon\to 0$; Costello--Gwilliam
Vol~I Ch.~3 Thm.~8.0.1 guarantees existence of a one-loop renormalised BV
theory with this counterterm.

\ClaimStatusTheorem\ \emph{Explicit coefficient of the running.} The
$\beta$-function-like coefficient for the wave-function renormalisation of
6D hCS at one loop is
\[
\boxed{\; Z_\cA^{(1)} \;=\; 1 \,-\, \hbar\,\frac{\Cas}{(4\pi)^3}\,\log(L/\epsilon)
\,+\,\cO(\hbar^2). \;}
\]
For $\gfrak=\sun$, $\Cas=2N$, giving coefficient $2N/(4\pi)^3 = N/(32\pi^3)$.

\subsection{Consistency with Vol III $\kch$ landscape}

$\Obs_{\hCS}$ on $\CC^3$ is the chain-level shadow of $\Phi$ at $d=3$ applied
to $A_{\CC^3}=\mathrm{Perf}(\CC^3)$, producing an $E_1$-chiral algebra
(Pattern 273 scope; not $E_2$ at $d\geq 3$). The Casimir coefficient $\Cas$
at one loop is the Kac--Moody-type central extension on the $E_1$-chiral
product, not a chiral-conformal weight. The abelian vanishing at all loops
witnesses the CY-D identification $\kch(A_{\CC^3})=\chi(\cO_{\CC^3})=1$
being additive with no Hodge-supertrace correction at $d=3$.

\subsection{References}
Costello~2011, \emph{Renormalization and effective field theory}, \S\S2.3,
9 (homotopy RG, heat-kernel regulator, BV master equation);
Costello--Gwilliam~2017 Vol~I, Ch.~3 (Bochner--Martinelli propagator,
Feynman rules for holomorphic field theories, one-loop existence theorem);
Costello~2013 \texttt{arXiv:1303.2632} \S\S9--10 (holomorphic-topological
Chern--Simons, Yangian from $\hbar$-deformation).
